\chapter{Conclusion et Perspectives}

\section{Rappel de la Problematique}

L'acces a l'information juridique marocaine en langue arabe reste un defi majeur. Les textes de loi sont disperses entre le Bulletin Officiel, les portails gouvernementaux, et les recueils prives. La majorite des documents historiques existe uniquement au format PDF scanne, rendant la recherche plein-texte impossible sans traitement prealable. Les solutions existantes --- legislation.ma, ADALA, portails du SGG --- reposent sur la recherche par mots-cles, sans capacite semantique ni generation automatisee de reponses.

Face a ce constat, le projet \textbf{Moroccan Legal RAG} a entrepris de construire le premier assistant juridique intelligent pour le droit marocain, fonde sur l'architecture RAG (Retrieval-Augmented Generation), capable de comprendre une question en arabe, de retrouver les articles de loi pertinents, et de generer une reponse juridique aneree dans les sources officielles.

\section{Synthese des Contributions}

Le projet a abouti a sept contributions principales :

\begin{enumerate}
    \item \textbf{Constitution d'un corpus juridique marocain structure en arabe.} Un pipeline complet d'extraction a ete developpe : PDF scanne $\to$ rendu PyMuPDF $\to$ OCR EasyOCR $\to$ correction par modele de vision (VLM Gemma 4) $\to$ nettoyage structurel (LegalTextCleaner) $\to$ normalisation arabe (ArabicNormalizer) $\to$ structuration en articles (LegalStructureParser). Le corpus final totalise \textbf{2\,357 articles} repartis sur 7 domaines juridiques, au format JSON avec metadonnees hierarchiques (chapitre, section, livre).

    \item \textbf{Benchmark systematique de 5 modeles d'embedding pour l'arabe juridique.} Evalue sur 65 variantes de requetes (arabe formel, darija, mots-cles, francais, anglais) contre 26 articles avec distracteurs difficiles. Le modele \textbf{E5-large} (1024 dimensions) a ete recommande avec P@1 = 72\%, MRR = 81\%, et Domain Accuracy = 89\%. Le modele BGE-M3 (P@1 = 69\%) constitue un second choix interchangeable ($p = 0.08$, non significatif). Un gap significatif de 29\% a ete observe entre l'arabe formel (91\%) et le darija (65\%), confirmant le defi de la diglossie.

    \item \textbf{Comparaison de 3 strategies de retrieval sur 32 questions juridiques.} Les strategies Dense (pgvector cosine) et RRF (Reciprocal Rank Fusion) ont ete evaluees sur le corpus complet. Resultats : Hit@5 = 65-69\%, Recall@5 = 57-59\%, MRR@5 = 0.46-0.57. La strategie BM25 basee sur PostgreSQL FTS a echoue (0 resultat sur 32 questions), conduisant au developpement d'un tokenizer arabe specialise avec 60+ stop words juridiques et integration de rank-bm25. Le re-benchmarking post-correctif est en cours.

    \item \textbf{Deploiement d'un LLM quantifie leger.} Le modele Gemma 4 2B a ete quantifie en precision Q4\_K\_XL au format GGUF, reduisant l'empreinte memoire de $\sim$4 Go (FP16) a $\sim$1.1 Go (reduction de 73\%) avec une perte de qualite inferieure a 2\%. L'inference est assuree par llama.cpp en mode CPU+GPU hybride, avec une latence de generation de $\sim$6 secondes par reponse sur GPU T4 Colab.

    \item \textbf{Service d'embedding GPU distant avec fallback local.} L'embedding E5-large est servi via Google Colab (GPU T4 gratuit), expose par Cloudflare Tunnel. En cas d'indisponibilite, le systeme bascule automatiquement sur un embedding CPU local (E5-base, 768 dimensions). Le debit mesure est de 42 textes/s en mode GPU contre 78 textes/s en mode CPU.

    \item \textbf{Benchmark RAG complet avec 5 metriques.} Le pipeline complet (retrieval + generation) a ete evalue sur 32 questions avec les metriques Context Precision, Context Recall, MRR, Answer Relevance, Faithfulness, et Answer Correctness. L'analyse a revele que la metrique de fidelite (token-overlap) sous-estime significativement la qualite des reponses, l'enrichissement contextuel du LLM (synthese multi-articles, citations, reformulation) etant un comportement souhaite mais non detecte par le chevauchement lexical.

    \item \textbf{Interfaces multiples et infrastructure conteneurisee.} Une API REST (FastAPI, 6 endpoints), une interface web (SvelteKit 5 + Tailwind CSS 4), une interface terminal (Textual), et des scripts CLI completes. L'ensemble est deployable via Docker Compose avec PostgreSQL 16 + pgvector 0.7.4.
\end{enumerate}

\section{Limites et Enseignements}

\subsection{Limites Identifiees}

Le Tableau~\ref{tab:limites} synthetise les limites identifiees lors du developpement et de l'evaluation du systeme :

\begin{table}[h]
\centering
\renewcommand{\arraystretch}{1.3}
\begin{tabular}{|p{3.5cm}|p{4cm}|p{4.5cm}|}
\hline
\rowcolor{LightBlue}
\textbf{\color{PrimaryBlue}Limite} & \textbf{\color{PrimaryBlue}Impact mesure} & \textbf{\color{PrimaryBlue}Piste d'amelioration} \\
\hline
Gap Darija--MSA & Baisse de P@1 de 29\% (91\% $\to$ 65\%) & Fine-tuning embedding sur corpus darija juridique \\
BM25 non operationnel & RRF = Dense effectivement & Re-benchmark apres correctif tokenizer arabe \\
Domaine fiscal inaccessible & 0\% sur toutes les metriques (tax) & Aligner taxonomie benchmark et base de donnees \\
LLM 2B parametres & Qualite de generation variable, Faithfulness $\sim$22\% (token-overlap) & Upgrade vers modele 7B+ (Gemma 7B, DeepSeek) \\
Latence tunnel Cloudflare & $\sim$5.5 s sur $\sim$12 s de pipeline total & Deploiement embedding local (CPU ou GPU) \\
Metrique de fidelite trompeuse & Sous-estimation systematique de la qualite & Implementer LLM-as-judge pour evaluation semantique \\
Corpus incomplet & Couverture partielle (7 domaines) & Extension continue (travail, urbanisme, sante) \\
\hline
\end{tabular}
\caption{Limites identifiees et pistes d'amelioration}
\label{tab:limites}
\end{table}

\subsection{Lecons Apprises}

Le projet a genere plusieurs enseignements methodologiques et techniques :

\begin{enumerate}
    \item \textbf{L'arabe juridique est un defi NLP majeur.} Meme les meilleurs embeddings multilingues (E5-large) peinent a discriminer des articles juridiquement proches mais distincts. La normalisation morphologique (alefs, tatweel, diacritiques) est indispensable mais insuffisante sans un tokenizer adapte au domaine.

    \item \textbf{La qualite OCR est le facteur limitant principal.} La correction par VLM (Gemma 4) ameliore significativement la qualite du texte, mais ne compense pas totalement les defauts des scans de basse qualite. L'investissement dans le pre-traitement des images (debruiter, redresser, binariser) est rentable.

    \item \textbf{La quantification GGUF est remarquablement efficace.} Le passage de FP16 a Q4\_K\_XL reduit l'empreinte memoire de 73\% avec une perte de qualite inferieure a 2\%, rendant viable le deploiement d'un LLM sur materiel modeste (GPU T4 gratuit, voire CPU).

    \item \textbf{Le RAG hybride est superieur a chaque strategie isolee.} La fusion Dense + BM25 via RRF combine la comprehension semantique et la precision terminologique, sans parametre a regler. L'echec du BM25 PostgreSQL souligne l'importance d'une tokenisation adaptee a la langue.

    \item \textbf{Colab + Cloudflare Tunnel est viable pour le prototypage, pas pour la production.} La latence de $\sim$5.5 s par appel d'embedding est acceptable pour l'experimentation mais inadaptee a un service interactif.

    \item \textbf{Un benchmark rigoureux est indispensable.} Les trois campagnes d'evaluation (embedding, strategies, RAG complet) ont permis de guider objectivement les choix techniques et d'identifier les faiblesses du systeme.
\end{enumerate}

\section{Comparaison avec l'Existant}

Le Tableau~\ref{tab:comparaison} positionne le systeme Moroccan Legal RAG par rapport aux alternatives existantes :

\begin{table}[h]
\centering
\renewcommand{\arraystretch}{1.3}
\begin{tabular}{|l|c|c|c|c|c|}
\hline
\rowcolor{LightBlue}
\textbf{\color{PrimaryBlue}Capacite} & \textbf{\color{PrimaryBlue}Google} & \textbf{\color{PrimaryBlue}legislation} & \textbf{\color{PrimaryBlue}ChatGPT } & \textbf{\color{PrimaryBlue}LLM } & \textbf{\color{PrimaryBlue}Notre systeme} \\
\hline
Recherche semantique arabe & Non & Non & Non & Non & \textbf{Oui} \\
Reponses generees & Non & Non & Oui & Oui & \textbf{Oui} \\
Citations d'articles & Non & Non & Parfois  & Parfois  & \textbf{Oui (verifiables)} \\
Ancrage dans les sources & -- & -- & Non & Non & \textbf{Oui} \\
Open-source & -- & -- & Non & Variable & \textbf{Oui} \\
Deployable localement & -- & -- & Non & Variable & \textbf{Oui (Docker)} \\
Corpus juridique marocain & Partiel & Oui & Non & Non & \textbf{Oui (7 domaines)} \\
\hline
\end{tabular}
\caption{Comparaison du systeme avec les alternatives existantes}
\label{tab:comparaison}
\end{table}

Les avantages distinctifs du systeme sont :
\begin{itemize}
    \item \textbf{Recherche semantique en arabe} : Comprehension du sens de la question, au-dela des mots-cles exacts.
    \item \textbf{Reponses tracables} : Chaque reponse cite les articles sources avec leur numero et la loi de reference.
    \item \textbf{Ancrage factuel} : La generation est contrainte par les documents retrouves, eliminant le risque d'invention d'articles.
    \item \textbf{Deploiement autonome} : Architecture 100\% conteneurisee, deployable sur infrastructure locale ou cloud.
\end{itemize}

\section{Perspectives}

Les perspectives de developpement sont organisees selon trois horizons temporels.

\subsection{Court terme (1-3 mois)}

\begin{itemize}
    \item \textbf{Re-benchmarking BM25} : Executer le benchmark RAG apres l'application du correctif de tokenisation arabe pour mesurer les performances reelles du canal BM25 et de la fusion RRF.

    \item \textbf{Correction du domaine fiscal} : Aligner la taxonomie des domaines entre le benchmark et la base de donnees (articles fiscaux actuellement stockes sous \texttt{domain = `civil'}).

    \item \textbf{Deploiement embedding local} : Remplacer le tunnel Cloudflare par un service d'embedding local (E5-base sur CPU, ou GPU local si disponible) pour reduire la latence de $\sim$5.5 s a $\sim$100 ms.

    \item \textbf{Amelioration du frontend} : Ajout de l'historique des conversations, des filtres avances (date, juridiction), et de l'export des reponses en PDF.
\end{itemize}

\subsection{Moyen terme (3-6 mois)}

\begin{itemize}
    \item \textbf{Fine-tuning embedding sur corpus juridique} : Specialiser E5-large ou BGE-M3 sur le corpus juridique marocain pour reduire le gap Darija--MSA et ameliorer la discrimination entre articles proches.

    \item \textbf{RAG multi-tours conversationnel} : Implementer la conservation du contexte de conversation pour permettre des questions de suivi et des clarifications progressives.

    \item \textbf{Evaluation humaine} : Constituer un panel de juristes marocains pour evaluer qualitativement les reponses generees sur les criteres de pertinence, exactitude juridique, et clarte.

    \item \textbf{Remplacement de la metrique de fidelite} : Substituer le chevauchement lexical par une evaluation LLM-as-judge capable de distinguer l'enrichissement contextuel legitime de l'hallucination.
\end{itemize}

\subsection{Long terme (6-12 mois)}

\begin{itemize}
    \item \textbf{Modele juridique specialise} : Fine-tuning d'un LLM (Gemma 7B ou equivalent) sur le corpus juridique marocain pour ameliorer la qualite de generation et la precision des references legales.

    \item \textbf{API publique} : Exposition du systeme via une API documentee avec SDK, destinee aux professionnels du droit (avocats, juristes d'entreprise, magistrats).

    \item \textbf{Mise a jour automatique} : Veille automatisee du Bulletin Officiel marocain avec detection des nouvelles lois, extraction, et re-indexation automatique du corpus.

    \item \textbf{Extension multi-pays} : Adaptation du systeme aux droits tunisien et algerien, partageant des racines juridiques communes avec le droit marocain (droit civil d'inspiration francaise, droit musulman).

    \item \textbf{Agent conversationnel juridique} : Evolution vers un agent capable de dialoguer avec l'utilisateur pour preciser sa situation juridique avant de fournir une reponse ciblee.
\end{itemize}

\section{Conclusion Generale}

Le projet Moroccan Legal RAG a demontre la faisabilite et la pertinence d'un assistant juridique intelligent pour le droit marocain en langue arabe. En combinant les avancees les plus recentes en NLP arabe (embedding multilingue E5-large), en recherche d'information (recherche hybride Dense + BM25 + RRF), et en generation de texte (LLM Gemma 4 quantifie), le systeme offre une solution operationnelle a un probleme concret : l'acces du citoyen et du professionnel a l'information juridique marocaine.

Les resultats quantitatifs sont encourageants :

\begin{itemize}
    \item \textbf{65-69\% de taux de succes} (Hit@5) pour la recherche d'articles pertinents sur le corpus complet de 2\,357 articles.
    \item \textbf{89\% de precision de domaine} : le systeme identifie correctement le domaine juridique dans la grande majorite des cas.
    \item \textbf{72\% de Precision@1} sur le benchmark d'embedding controle (E5-large).
    \item \textbf{Deploiement sur materiel modeste} : $\sim$1.1 Go pour le LLM, $\sim$2 Go de RAM pour PostgreSQL, conteneurisation Docker complete.
\end{itemize}

Les limites identifiees --- gap dialectal, latence du tunnel Cloudflare, couverture partielle du corpus --- definissent une feuille de route claire pour les ameliorations futures. La correction du BM25, le deploiement local de l'embedding, et l'extension du corpus constituent les priorites immediates.

Au-dela de ses resultats techniques, ce projet illustre le potentiel de l'architecture RAG pour democratiser l'acces au droit dans les pays ou la barriere linguistique et la dispersion des sources constituent un obstacle majeur. La methodologie developpee --- extraction OCR+VLM, normalisation arabe, embedding multilingue, recherche hybride, generation aneree --- est transposable a d'autres corpus juridiques arabophones (Tunisie, Algerie, Egypte) et, plus largement, a tout domaine ou l'information textuelle est abondante mais structurellement inaccessible.